Два человека упали с крыши пятиэтажного дома, новостройки. Кажется, школы. Они съехали по крыше в сидячем положении до самой кромки и тут начали падать.

Их падение раньше всех заметила Ида Марковна. Она стояла у окна в противоположном доме и сморкалась в стакан. И вдруг она увидела, что кто-то с крыши противоположного дома начинает падать. Вглядевшись, Ида Марковна увидела, что это начинают падать сразу целых двое. Сове>

В это время падающих с крыш увидела другая особа, живущая в том же доме, что и Ида Марковна, но только двумя этажами ниже. Особу эту тоже звали Ида Марковна. Она, как раз в это время, сидела с ногами на подоконнике и пришивала к своей туфле пуговку. Взглянув в окно, она>

Схватив клещи, Ида Марковна опять подбежала к окну и выдернула гвоздь. Теперь окно легко распахнулось. Ида Марковна высунулась из окна и увидела, как падающие с крыши со свистом подлетали к земле.

На улице собралась уже небольшая толпа. Уже раздавались свистки, и к месту ожидаемого происшествия не спеша подходил маленького роста милиционер. Носатый дворник суетился, расталкивая людей и поясняя, что падающие с крыши могут ударить собравшихся по головам.

К этому времени уже обе Иды Марковны, одна в платье, а другая голая, высунувшись в окно, визжали и били ногами.

И вот наконец, расставив руки и выпучив глаза, падающие с крыши ударились об землю.

Так и мы иногда, упадая с высот достигнутых, ударяемся об унылую клеть нашей будущности.

\begin{flushright}
7 сентября 1940~г.
\end{flushright}

